\subsection{Mapped CON-bin columns} \label{sec:conbin-mmap}

The mapping is \texttt{MAP\_PRIVATE}.
The advice calls are \texttt{madvise}.
\texttt{MADV\_} and \texttt{MAP\_} are the prefixes on those constants.
Memory advice (MADV) names the \texttt{madvise} requests.
The mapping flag (MAP) names the \texttt{mmap} flags.
Non-uniform memory access (NUMA) names the machine layout.
Tables~\ref{tab:conbin-mmap-sum} and \ref{tab:conbin-scale} report the
resident set size (RSS) as \texttt{VmHWM}.
\texttt{VmHWM} names the high-water mark of the resident set for each
process.
A CON-bin file begins with one host page.
On this machine that page is 4096 bytes.
The first 40 bytes hold the magic, the version, the flags, the frame
count, the atom count and the payload offset.
Positions of all 100 frames follow as one little-endian binary64 run,
then velocities.
The payload is
\(100 \times 10^{6} \times 3 \times 8 \times 2 = 4\,800\,000\,000\)
bytes.
\texttt{read\_conbin} returns those two runs as writable NumPy arrays
over a \texttt{MAP\_PRIVATE} mapping.
A store copies the touched page and leaves the file unchanged.
The map is the NumPy base of both arrays and is unmapped when both
arrays are freed.
\texttt{mmap(2)} leaves visibility of a later write by another process
up to the implementation, for a page the caller has not stored to.
Truncation of the file can raise \texttt{SIGBUS}.
\texttt{read\_conbin\_copy} fills a caller-owned buffer with
\texttt{pread}.
h5py reads two contiguous datasets.
The raw chunk cache stays unset, because those datasets are not chunked.

Mapping advice lists \texttt{MADV\_SEQUENTIAL}, \texttt{MADV\_HUGEPAGE},
\texttt{MADV\_WILLNEED}, then \texttt{MADV\_POPULATE\_READ}.
HUGEPAGE applies when transparent large pages are in madvise mode.
File-backed large pages also require
\texttt{CONFIG\_READ\_ONLY\_THP\_FOR\_FS}.
That option is unset, and
\texttt{FilePmdMapped} stayed 0.
The population call still faults base pages.

Those base pages are what Table~\ref{tab:conbin-mmap-nosum} times, each
reader in a fresh process after \texttt{drop\_caches} 3.
That drop clears the page cache.
The host cache was not controlled.
The warm sample in that table is the next read in that same process.
There is no untimed warm-up.
Repeat 0 is the first touch of that cell and is kept.
Six repeats.
The reported value is the upper median, sorted index 3.
The rate numerator is the \(4\,800\,000\,000\)-byte payload for every
reader.
Table~\ref{tab:conbin-mmap-nosum} stops when the arrays exist.
The following table keeps a float64 sum of both columns inside the
clock.
Every sum in Table~\ref{tab:conbin-mmap-nosum} was
\(-1.12499999625 \times 10^{16}\).

The memory-mapped reads use the same Xeon.
L3 is 320 MiB.
The page size is 4096 bytes and \texttt{Hugepagesize} is 2048 kB.
rustc is 1.88.0, Python is 3.12.3, h5py is 3.16.0 and NumPy is 1.26.4.

\begin{table}[ht] \centering

\caption{CON-bin and h5py, clock stopped when the arrays exist.
Upper median of six repeats, range in parentheses.
h5py in this table is one reader.
RSS is the resident set after that read, in kB.
Cold means the page cache was dropped before the read. The host cache was not controlled.}
\label{tab:conbin-mmap-nosum}
\resizebox{\linewidth}{!}{%
\begin{tabular}{@{}llrrrr@{}}
\toprule
Reader & Threads & Cache & Upper median (s) & MB/s & RSS \\
\midrule
pread & 1 & cold & 3.0708 (2.8267 to 3.8110) & 1563.1 & 4735008 \\
pread & 1 & warm & 2.4222 (2.3441 to 2.6418) & 1981.7 & 4735012 \\
mapped & 1 & cold & 1.8443 (1.7338 to 1.9078) & 2602.6 & 4734900 \\
mapped & 1 & warm & 0.2415 (0.2307 to 0.2747) & 19871.9 & 4734904 \\
h5py & 1 & cold & 3.0729 (2.6581 to 4.2668) & 1562.0 & 4735788 \\
h5py & 1 & warm & 2.4124 (2.3719 to 2.5334) & 1989.8 & 4736300 \\
pread & 8 & cold & 1.5140 (1.2922 to 1.8403) & 3170.4 & 4735264 \\
pread & 8 & warm & 0.4114 (0.3843 to 0.4197) & 11668.2 & 4735264 \\
mapped & 8 & cold & 1.5664 (1.3673 to 1.9114) & 3064.4 & 4735360 \\
mapped & 8 & warm & 0.1101 (0.1072 to 0.1165) & 43582.0 & 4735360 \\
\bottomrule
\end{tabular}}
\end{table}
% Periods in the table are decimals.

\begin{table}[ht]
\centering
\caption{pread, plain map, and h5py from the same session as
Table~\ref{tab:conbin-scale}.
The clock includes a float64 sum of both columns.
Upper median of six repeats.
Warm is a new process.
RSS is \texttt{VmHWM}, in kB. h5py in this table is the 1-thread cell of
that session.
Cold means the page cache was dropped before the read. The host cache was not controlled.}
\label{tab:conbin-mmap-sum}
\resizebox{\linewidth}{!}{%
\begin{tabular}{@{}llrrrr@{}}
\toprule
Reader & Threads & Cache & Upper median (s) & MB/s & RSS \\
\midrule
pread & 1 & cold & 3.2126 (2.8225 to 3.8666) & 1494.1 & 4689984 \\
pread & 1 & warm & 2.9315 (2.7362 to 3.0299) & 1637.4 & 4690000 \\
mapped & 1 & cold & 2.3343 (2.1735 to 2.5727) & 2056.3 & 4690076 \\
mapped & 1 & warm & 0.9173 (0.9140 to 0.9887) & 5232.5 & 4690016 \\
h5py & 1 & cold & 3.3113 (3.0086 to 7.0523) & 1449.6 & 4736628 \\
h5py & 1 & warm & 2.7408 (2.6431 to 2.8928) & 1751.3 & 4736684 \\
pread & 8 & cold & 2.0895 (2.0158 to 3.1192) & 2297.2 & 4690432 \\
pread & 8 & warm & 1.1063 (1.0692 to 1.2108) & 4338.6 & 4690436 \\
mapped & 8 & cold & 1.9069 (1.8371 to 2.0154) & 2517.2 & 4690408 \\
mapped & 8 & warm & 0.9293 (0.9121 to 0.9431) & 5165.2 & 4690424 \\
\bottomrule
\end{tabular}}
\end{table}
% Periods in the table are decimals.

\subsection{Threaded readahead}
\label{sec:conbin-scale}

\texttt{max\_sectors\_kb} is 1280.
Each hardware queue has depth 256 (\texttt{nr\_requests}).

This host has one hardware queue, because \texttt{nr\_hw\_queues} is absent.
The default is one stream.
Workers fault a piece with \texttt{MADV\_POPULATE\_READ}, then sum it.
\texttt{mincore} only decides whether the whole file is already
resident.
The piece sums are added in file order.
This disk records the rotational flag as 1.

\texttt{read\_conbin} returns writable arrays over a
\texttt{MAP\_PRIVATE} mapping (\texttt{memmap2} \texttt{map\_copy}).
A store copies the page.
The file bytes stay as written.
The map is the NumPy base of both arrays and is unmapped when both are
freed.
Advice \texttt{MADV\_HUGEPAGE} is requested.
\texttt{CONFIG\_READ\_ONLY\_THP\_FOR\_FS} is unset, which is the switch required before
khugepaged can collapse file pages.
\texttt{FilePmdMapped} stays 0.

The plain map is one \texttt{MADV\_POPULATE\_READ} per worker.
\texttt{O\_DIRECT} is the lazy two-buffer read.
h5py 3.16.0 reads two contiguous datasets, then \texttt{numpy.sum}
(NumPy 1.26.4).
The datasets are not chunked, so the raw chunk cache is left unset.
\texttt{H5is\_library\_threadsafe} returns 0 for this build, so each
worker is a process with its own file handle.
The readahead sum and the h5py sum are
\(-1.12499999625 \times 10^{16}\).
The plain map and \texttt{O\_DIRECT} add one running total and reach
\(-1.1249999941777216 \times 10^{16}\).
The last digits differ because floating-point addition depends on order.
rustc is 1.88.0 and Python is 3.12.3.
The clock includes the sum.
Six repeats.
The reported sample is the upper median, sorted index 3.
The numerator is \(4800000000\) bytes.
Every reader in the scale table runs in its own cgroup.
The RSS column is \texttt{VmHWM} of the reader, in kB.
For an h5py pool it is the parent's mark plus the sum of the workers'
marks, an upper bound on the concurrent peak.
\texttt{VmHWM} counts the pages a process maps.
Page cache that no process maps stays outside that mark.
The h5py cold read fills the page cache and keeps its own arrays.
At 1 thread that cgroup high-water mark reaches 9431144 kB.
The mapped read of the same file reaches 4709088 kB.
The JSON records that cgroup mark for every sample.
The probe's queue and its cache type both read write through.
The cache probe read 1 GiB of the benchmark file at 397.7 MB/s after
\texttt{drop\_caches} 3.
After an \texttt{O\_DIRECT} write and the same drop, the buffered read
was 2934.1 MB/s.
The two later \texttt{O\_DIRECT} reads of that file were 2188.4
MB/s and 2246.7 MB/s.
The other two files' \texttt{O\_DIRECT} pairs were 5924.6
MB/s then 2271.3 MB/s, and 2187.7 MB/s then 2158.6 MB/s.
Cold means the page cache was dropped before the read. The host cache was not controlled.
Host RAM is unknown.
A file written and never read, then read by one stream after that drop,
took 1.831 s, then 1.758 s, then 1.609 s.
A file read many times took 1.605 s and 1.708 s.
Each recorded \texttt{read\_bytes} of the whole file.
\texttt{Cached} was 312000 kB before each of those reads.
The never-read file's first read was the slowest of the five, 1.14 times
the fastest.
The page cache was dropped before every cold sample.
The first read of each file in this session is the slowest sample of its
cell.
h5py at 1 thread took 7.0523 s (680.6 MB/s), 2.13 times its upper
median.
One CON-bin stream at 6 threads took 7.2233 s (664.5 MB/s), 7.11 times
its upper median.
The next is O\_DIRECT at 2 threads, repeat 4, at 1.54 times.
No sample records whether a host cache served it.
Repeat 0 of the other cold h5py cells is 0.92 to 1.19 times its
upper median.
The samples are kept.
The upper median absorbs the first reads.
\texttt{fio} 3.36 used \texttt{io\_uring}, \texttt{direct=1},
\texttt{rw=read} and a 1310720-byte block.
One job at depth 1 read at 1402.6 MB/s.
One job at depth 256 read at 2680.0 MB/s.
Eight jobs at depth 256, on disjoint spans, read at 3309.4 MB/s.
One stream took 1.0515 s in the cold column at 1 thread (4564.8
MB/s) and 1.0217 s at 8 threads (4698.2 MB/s).
The slowest one-stream rate in that column is 4424.9 MB/s, at 5 threads.
A stream per span took 1.3572 s at 1 thread and 1.2101 s at 8 threads.
The probe read the benchmark file at 397.7 MB/s after the same drop.
The cold rates above 4000 MB/s in Table~\ref{tab:conbin-scale}
therefore came from a host cache.
They are not device rates.
The default stays one stream, the simpler plan.
O\_DIRECT still meets a host cache.
The measured rate is a cache rate.

\begin{table}[ht] \centering
\caption{Plain map, one stream, a stream per span, \texttt{O\_DIRECT}, and h5py.
Modes are shuffled inside each repeat.
Warm is a new process.
The parent waits on a barrier after the h5py workers import, then starts
the clock.
Each h5py worker reads one slice.
The resident set size (RSS) is \texttt{VmHWM}, in kB.
An h5py pool adds the parent's mark and each worker's mark.
The h5py 1-thread cell is the same cell as in
Table~\ref{tab:conbin-mmap-sum}.
Cold means the page cache was dropped before the read. The host cache was not controlled.}
\label{tab:conbin-scale}
\resizebox{\linewidth}{!}{%
\begin{tabular}{@{}llrrrr@{}}
\toprule
Reader & Threads & Cache & Upper median (s) & MB/s & RSS \\
\midrule
plain map & 1 & cold & 2.3343 (2.1735 to 2.5727) & 2056.3 & 4690076 \\
readahead & 1 & cold & 1.0515 (0.9978 to 1.2141) & 4564.8 & 4690416 \\
split & 1 & cold & 1.3572 (1.2174 to 1.4151) & 3536.7 & 4690240 \\
O\_DIRECT & 1 & cold & 1.6519 (1.4851 to 1.9446) & 2905.7 & 4690128 \\
h5py & 1 & cold & 3.3113 (3.0086 to 7.0523) & 1449.6 & 4736628 \\
plain map & 1 & warm & 0.9173 (0.9140 to 0.9887) & 5232.5 & 4690016 \\
readahead & 1 & warm & 0.8734 (0.8521 to 0.9273) & 5495.9 & 4690272 \\
split & 1 & warm & 0.9000 (0.8694 to 1.0086) & 5333.5 & 4690248 \\
O\_DIRECT & 1 & warm & 1.7614 (1.4944 to 1.8873) & 2725.1 & 4690136 \\
h5py & 1 & warm & 2.7408 (2.6431 to 2.8928) & 1751.3 & 4736684 \\
plain map & 8 & cold & 1.9069 (1.8371 to 2.0154) & 2517.2 & 4690408 \\
readahead & 8 & cold & 1.0217 (0.9683 to 1.1867) & 4698.2 & 4690812 \\
split & 8 & cold & 1.2101 (0.9608 to 1.4188) & 3966.5 & 4689924 \\
O\_DIRECT & 8 & cold & 1.5320 (1.4677 to 1.5830) & 3133.2 & 4690436 \\
h5py & 8 & cold & 1.4268 (1.2302 to 1.8428) & 3364.2 & 5124632 \\
plain map & 8 & warm & 0.9293 (0.9121 to 0.9431) & 5165.2 & 4690424 \\
readahead & 8 & warm & 0.1779 (0.1726 to 0.1870) & 26977.5 & 4690248 \\
split & 8 & warm & 0.1821 (0.1733 to 0.1927) & 26363.2 & 4690372 \\
O\_DIRECT & 8 & warm & 1.5250 (1.4656 to 1.6083) & 3147.6 & 4690436 \\
h5py & 8 & warm & 0.3989 (0.3914 to 0.4091) & 12033.0 & 5126924 \\
\bottomrule
\end{tabular}}
\end{table}
% Periods in the table are decimals.
